\section{A taxonomy of persistent state}\label{sec:statetax}


\regime{Substance}{the kinds of persistent state an always-on agent holds, and which axes each stresses.}
Persistent-state taxonomy starts with a harder question than memory taxonomy: what counts as state? A memory survey can take its object for granted: stored information that can be retrieved. A persistent-state survey cannot, because the records that govern an always-on agent's next action are wider than those it would call memories. A revoked API token, a half-completed purchase, a lapsed consent, a procedural skill learned in an environment that has since changed, and a calendar trigger about to fire all condition future behavior, yet none is naturally described as a remembered fact. The taxonomy here enumerates the persistent-state types an always-on agent maintains, grounding each in representative work, naming the failure mode it is prone to, and identifying which of the six state axes (authority, scope, mutability, provenance, recoverability, actionability) it stresses most. Many of these types descend from older cognitive-science categories; the novelty is the obligation to govern them as one interacting state rather than as separable memory forms. The governance burden is uneven: the types the field studies most are the ones whose failures are most recoverable, and the types it studies least are the ones whose failures are least.

We use the three-layer partition summarized in Table~\ref{tab:state-tax}. \emph{Object state} is the thing the agent carries forward: episodic traces, semantic facts, preferences and user models, procedural skills, workflow and task commitments, shared state, and trigger conditions. The \emph{governance envelope} is the metadata and capability surface that travels with those objects: authority, scope, credentials, provenance, retention, deletion, rollback handles, trust, and temporal validity. \emph{External commitments} are externally committed effects linked to internal state: messages, payments, files, database rows, tickets, and other consequences the agent does not fully own once committed. For exposition, the subsections below discuss recurring categories, while Table~\ref{tab:state-tax} further decomposes the governance envelope into its main fields. This split matters because several items that a flat taxonomy would list as state types, such as policy, permission, provenance, and deletion handles, are better modeled as envelope fields over every state object, while side effects are not internal state at all but external consequences that must remain traceable to it. Throughout, we model a persistent-state item as the record $s = \langle v, a, c, m, p, r, k, \tau \rangle$ introduced in the foundations section: a value $v$ under an authority $a$, a scope $c$, a mutability class $m$, a provenance chain $p$, a recoverability handle $r$, an actionability type $k$, and a logical timestamp $\tau$. A state type is then a region of this record space that shares an actionability type and a characteristic profile over the other axes, and a failure mode is the violation of an invariant that region is most exposed to.

\subsection{Memory types: the classical core}

\subsubsection{Episodic traces}

Episodic traces are the raw record of what happened: conversation turns, tool calls and their outputs, environment observations, and the timestamps that order them. They are the most literal sense of memory and the foundation the other types derive from, since a semantic fact is usually an abstraction over episodes and a procedural skill a distillation of successful ones. The category descends directly from Tulving's separation of episodic memory, indexed by time and context, from context-free semantic memory \citep{tulving1972episodic}, a distinction the Soar architecture later operationalized by giving an agent distinct episodic and semantic stores within a single decision cycle \citep{laird2012soar}, and that the working-memory tradition refined through the episodic buffer that binds multimodal experience into integrated, retrievable episodes \citep{baddeley2000episodic}. Modern agent frameworks inherit this structure: CoALA organizes language agents around explicit memory modules whose episodic component logs interaction history for later retrieval \citep{sumers2024cognitivearchitectureslanguageagents}, and broad agent surveys treat the episodic log as the substrate from which all higher memory is consolidated \citep{xi2023riseandpotential}.

The characteristic failure of episodic traces is not forgetting but the opposite: undisciplined accumulation. An episodic store grows linearly with interaction, accumulates noise, can leak private content captured incidentally during ordinary turns, and loses the handles needed to act on or audit a trace later. Cognitive literature already suggests the response: raw episodes should be selectively consolidated rather than retained verbatim. Complementary-learning-systems theory argues that a fast episodic store must interleave with slow semantic consolidation to avoid catastrophic interference \citep{mcclelland1995complementary,kumaran2016complementary}, and the demonstration of offline hippocampal replay during sleep \citep{wilson1994reactivation} supplies the biological mechanism that agent experience-replay schemes imitate. Episodic traces stress recoverability: an episode is only useful downstream if it retains the provenance and retrieval handles that let a later decision be traced back to it, yet the typical agent log stores surface text without the source, timestamp, and trust tier that the provenance chain $p$ and recoverability handle $r$ require. The open problem is that the field captures episodes abundantly but rarely as governable records; the trace is written without the metadata any later forget, audit, or rollback operation would need, so an episodic store has plenty of $v$ and little else.

\subsubsection{Semantic facts}

Semantic facts are the context-free assertions an agent extracts and retains about its users, its projects, and the world: that a user lives in a particular city, that a project uses a particular framework, that an account is in a particular state. They are the type a personalization system most often means by long-term memory, and the recent personalization literature has made them the center of evaluation. PersonaMem tracks an evolving user persona over many interactions and scores whether the agent recalls the current profile \citep{jiang2025personamem}; PrefEval tests whether a model follows a stated fact or preference across long context \citep{zhao2025prefeval}; LaMP established the earlier baseline of profile-conditioned generation \citep{salemi-etal-2024-lamp}; and production-scale systems such as the industrial long-term semantic memory of \citet{xu2026linkedinmemory} show that distilling durable user facts from noisy longitudinal signals is feasible at deployment scale. CarMem narrows the type usefully by bounding semantic facts to declared categories, letting it sharply cut redundant and contradictory stored preferences through scheduled maintenance \citep{kirmayr2025carmem}, and surveys of personalized LLM agents codify the requirement that such facts persist coherently across sessions \citep{xu2026personalizedllmpoweredagentsfoundations}.

The failure modes of semantic facts are staleness, contradiction, and authority confusion, and they sharpen as the timescale lengthens. A fact that was true last month may be false today, and the agent that recalls it with high fidelity is then confidently wrong; DynamicMem shows that user attributes, habits, and preferences drift on different timelines driven by external events, so a single decay rate cannot keep a semantic store current \citep{xie2026dynamicmem}. The deeper problem the personalization literature surfaces, and the one most aligned with the thesis, is authority confusion: a benchmark that scores recall rewards an agent for retrieving a fact regardless of whether it still carries a current, unrevoked authority to influence action. MemProbe reframes the evaluation question by treating long-term memory as a post-interaction artifact to be audited for hidden user-state rather than judged only by downstream task success \citep{ma2026memprobe}, beginning to expose what recall metrics hide. Semantic facts stress mutability and authority: a fact must be revisable and supersedable as the world changes, and it must lose its license to act when the authority behind it lapses. The gap is that the field measures whether a fact can be recalled, not whether recalling it is still authorized; an agent can top a preference-following benchmark and still have no mechanism to revoke a fact the user has retracted, an authority-monotonicity failure that recall benchmarks have no field to score.

\subsubsection{Preferences}

Preferences are a distinguished sub-class of semantic state: durable assertions about how the user wants the agent to behave, including style, habits, constraints, and priorities. We separate them from plain facts because their actionability type differs. A fact is evidence the agent reasons over, whereas a preference is a standing instruction the agent applies, so it more directly shapes action and raises the scope and authority stakes. The evaluation literature treats preference handling as a first-class capability: PrefEval probes whether a model infers and adheres to stated preferences over long multi-session context \citep{zhao2025prefeval}, MPT tests whether an agent reuses persisted preferences to fill under-specified tool-call arguments across sessions, decomposing the capability into recall, induction, and transfer \citep{yoon2026mpt}, and PAL-Bench evaluates assistants on subjective user traits accumulated over long-term interaction \citep{huang2025mempal}. POLAR shows the upside when the type is handled well, accumulating personalized context to resolve implicitly specified targets that a stateless agent could not \citep{lee2026polar}, and the foundations survey of personalized agents frames preference persistence as a core requirement rather than a feature \citep{xu2026personalizedllmpoweredagentsfoundations}.

The failure modes of preferences are over-personalization, scope errors, and drift, and the most dangerous is a security concern rather than a quality one. \citet{xu2026toxicchats} formalize unintended long-term state poisoning, in which routine interactions gradually weaken an agent's confirmation boundaries and silently expand the action scope a stored preference may govern: a preference for fast checkout, accreted over many turns, quietly becomes a license to skip confirmation on a purchase. This is the scope axis failing under accumulation, the scope-non-expansion invariant violated one harmless-looking turn at a time. Preferences therefore stress scope and actionability above all: the question is what the user wants, which tasks, tools, and time windows that preference may act on, and whether it is advisory or has hardened into an executable default. The gap is that preference benchmarks score adherence, treating stronger following as strictly better, while the always-on risk is over-adherence: a preference whose scope has crept or whose authority has lapsed should stop governing action, and no current preference benchmark scores the agent that correctly declines to apply a preference it once held.

\subsubsection{Procedural skills}

Procedural skills are reusable, executable competence: plans, scripts, trajectories, and reflections that an agent stores after success and replays later. The category again has cognitive roots, in the procedural-versus-declarative distinction that ACT-R formalizes through utility-weighted production rules competing with activation-decayed declarative chunks \citep{anderson2004actr}, and Soar's separate procedural store \citep{laird2012soar}. In LLM agents the type was popularized by Voyager, which discovers and banks reusable skills to drive open-ended embodied learning with transfer across tasks \citep{wang2023voyageropenendedembodiedagent}, and by reflection learners such as Reflexion, which convert failed trajectories into verbal lessons stored as persistent guidance \citep{shinn2023reflexionlanguageagentsverbal}. Contextual experience replay folds past trajectories back into context as in-context skill rather than externalizing them \citep{liu2025contextualexperiencereplay}, and MAGE manages procedural state as a hierarchical execution tree with explicit grow, compress, maintain, and revise operations \citep{chen2026mage}. Procedural skills carry the actionability axis to its extreme: a recalled fact is evidence the model may weigh, but a recalled skill is a callable action whose misuse has tool consequences, so the bar for authorizing, scoping, and reversing it is correspondingly higher.

The failure modes are procedure poisoning, unsafe reuse, and overfitting, united by the absence of a validation gate between storing a procedure and trusting it. A reflection induced from a single failed run can become a plausible but unproven durable rule, and a skill that succeeded in one environment can be replayed after the goal or environment has shifted, turning a correct past execution into a wrong present one. The corpus has begun to name these directly: skill drift is formalized as a contract violation, with a check on whether stored procedural skills still produce their role-bearing outcomes in an evolving environment \citep{fan2026skilldriftcontract}, MemP makes deprecation of stale procedures explicit, which is rare \citep{fang2026mempexploringagentprocedural}, and skill containment promotes a skill manifest into a mechanically checkable capability-containment proof via abstract interpretation and bounded model checking \citep{metere2026skillcontainment}. These are the closest the literature comes to governing the type, and they remain exceptions. Procedural skills stress authority, mutability, and recoverability jointly: a skill must be tested and scoped before reuse, deprecated when superseded, and reverted when it causes a regression. The gap is that skill libraries are optimized for cheap accumulation and reuse, not for testing before reuse or rollback after harm, so competence and risk accumulate together, and outside MemP and the drift-and-containment line almost no system re-validates a stored procedure at the moment of reuse.

\subsection{Control-relevant state: beyond memory}

The four types above are where a memory survey would often stop. The remaining categories make an always-on agent a stateful system rather than a memory-augmented one: they include live commitments, shared state, triggers, tool capabilities, authority, provenance, and irreversible external effects. Some are object states in their own right; others are envelope fields attached to every object state. This distinction is the survey's main taxonomic correction: governance is not one more memory type but the layer that determines whether any remembered object may act.

\subsubsection{Workflow and task-ledger state}

Workflow and task-ledger state is the agent's record of open commitments: tasks in progress, deadlines, dependencies, partially completed multi-step operations, and the obligations they imply. It lets an always-on agent resume a long-running goal after a restart or interruption rather than begin afresh, and it is invisible to a memory-form taxonomy because it is live control state, not remembered information. AppWorld stresses the type directly, requiring agents to track world-state changes across sessions and apply stored behavioral patterns to new contexts over stateful applications \citep{trivedi-etal-2024-appworld}, while MAGE's execution tree is in part a task ledger, deriving state from action trajectories so a long operation can be inspected and revised mid-flight \citep{chen2026mage}, and agent workflow memory induces reusable high-level procedures from past task structure \citep{wang2024agentworkflowmemory}. The transactional view is the most governance-aware: SagaLLM imports the database Saga model of compensable execution into multi-agent planning, so a multi-step workflow can be compensated rather than left half-applied when a step fails \citep{chang2025sagallm}.

The characteristic failures are false continuity and stale commitments: the agent believes a task is open that has been completed or cancelled, or it resumes an obligation whose preconditions have since changed. These are mutability and recoverability failures, because a task ledger must update atomically as the world changes and must support compensation when a partially applied workflow has to be unwound. The gap is that most agent memory work treats the ledger as ordinary recall, scoring whether the agent remembers the task rather than whether it correctly maintains, supersedes, and compensates the commitment; outside the transactional line of SagaLLM, the corpus rarely tests whether an interrupted workflow is resumed under still-valid preconditions or whether a cancelled obligation actually stops driving action.

\subsubsection{Policy and permission state}

Policy and permission state is the record of what the agent is allowed to do: user consent, access rights, role assignments, and the safety constraints that bound action. It is the purest carrier of the authority axis and one of the least studied state types in the corpus: authority appears in only $72$ of $435$ works. The works that address it directly are recent and few. Intent-driven authorization treats a user's expressed intent as a monotone, session-scoped policy that can only reduce, never expand, the credential authority granted to tool calls, a direct instantiation of the authority-monotonicity invariant \citep{zhu2026intentgovernedauth}. CaMeL enforces permission at the system level by separating control flow from data flow, so that what a tool may do does not depend on what untrusted text requests \citep{debenedetti2025camel}. SAMEP shares persistent context across sessions under cryptographic access control, raising the question of how permission travels with state across boundaries \citep{masoor2025samep}. These works establish that permission is enforceable but treat it as a static gate.

The failure modes are privilege drift and cross-user leakage: a permission granted for one purpose silently broadens to cover another, or one user's access rights bleed into another user's context through shared state. Both violate the authority and scope invariants, and both disappear in episodic agents only because permissions reset between tasks. The gap is severe and structural: the corpus has mechanisms to grant and check permission but almost none to revoke it dynamically, age it out, or audit whether an action was licensed by an authority current at the time. An agent can hold a permission long after the consent behind it has lapsed, and in the common case no mechanism notices, the authority-monotonicity invariant failing silently over time.

\subsubsection{Provenance and audit state}

Provenance and audit state is the metadata that makes every other type accountable: the sources, timestamps, lineage, transformations, and confidence attached to a record, plus the trail of who or what mutated it. It most directly enables the return arc of the lifecycle, because without provenance there is no basis for attributing a fact, verifying a claim, or tracing an action back to the records that justified it. A small but sharpening body of work targets it. MemLineage attaches a weighted derivation DAG and a Merkle log to agent memory and enforces an untrusted-path-persistence invariant that refuses actions descending from external sources \citep{ouyang2026memlineage}; immutable-memory designs use blockchain-indexed append-only ledgers with privilege lattices to make revisions protocol-blocked and historied \citep{wright2025immutablemem}; Causal Past Logic defines runtime-verifiable guards over causally visible stored state in distributed workflows \citep{bollig2026causalpast}; and MemProbe argues that provenance and hidden state should be audited as a first-class artifact \citep{ma2026memprobe}. The failure this type guards against is provenance-role collapse, in which flat-text memory loses source attribution, so the agent hallucinates authority and misattributes facts to the wrong source \citep{jin2026provenancerolecollapse}.

Provenance and audit state stress provenance preservation and, through it, recoverability: consolidation may compress a value, but it must retain enough lineage to verify, attribute, and later revise or revert the result. The gap is twofold. First, provenance is usually the first casualty of consolidation, because summarization collapses many sourced episodes into one unsourced summary, the corroboration-laundering failure the immutable-memory and lineage lines try to prevent. Second, immutability and deletion are in direct tension: an append-only provenance ledger that makes the past tamper-evident also makes a genuine forget request hard to honor, so a design that maximizes auditability can violate deletion propagation, and the corpus has not resolved how to reconcile the two. The escalation this type adds is that provenance is not one governance axis among several but the precondition for the others: without a preserved derivation chain there is nothing to delete against, nothing to attribute an action to, and nothing to roll back along, so the thinness here is what makes the thinness everywhere downstream irreparable.

\subsubsection{Shared and social memory}

Shared and social memory is persistent state that more than one party can read or write: group threads, team knowledge, role assignments, and the social context of multi-party interaction. It changes the governance problem qualitatively, because a fact written by one agent or user may act for another, so scope and authority must be enforced across principals rather than within one. The foundational demonstration is Generative Agents, whose observation, reflection, and retrieval memory in a shared environment produces believable social behavior, but with no access control, no authority framework, and unaudited memory open to injection \citep{park2023generativeagentsinteractivesimulacra}. MetaGPT's publish-subscribe message pool is a de-facto shared state, untyped and unaudited, with no deletion or revocation semantics \citep{hong2023metagpt}. The governance-aware response is recent: Collaborative Memory adds dynamic, field-level, role-anchored access control to multi-user memory \citep{rezazadeh2025collaborativememorymultiusermemory}, and governed shared memory enforces explicit authority and scope on multi-agent reads and writes while preserving read/write invariants \citep{margalit2026governedsharedmemorymultiagent}. The distinctive failure mode is propagation: memory contagion shows that evaluator bias stored in shared trajectories propagates cross-temporally to future agents, with no safe contamination threshold even under oracle consolidation \citep{liu2026memorycontagion}, and poison-propagation studies map how a single compromised entry spreads through agent-society topologies \citep{men2024troublemaker}.

Shared and social memory stresses scope and provenance jointly, with deletion propagation as the acute open problem: when one principal leaves a group or retracts a contribution, the deletion must reach every derived copy across every agent that read it, and shared-memory work rarely scores revocation or cross-thread leakage. In our coded corpus, the trajectory has moved from ungoverned shared substrates toward access-controlled ones, but deletion-propagation and authority-revocation guarantees across principals remain thin; this is the multi-principal form of the same invariants the survey names.

\subsubsection{Trigger state}

Trigger state is the record of when an always-on agent should act on its own: availability windows, proactive opportunities, and the conditions under which it should interrupt the user or fire a scheduled action. It most distinguishes an always-on agent from a reactive one, because a reactive agent acts only when prompted whereas an always-on agent must decide for itself when latent state warrants action, and it is barely studied. TriggerBench measures prospective memory directly, testing whether an agent spontaneously recalls and acts on a latent stored constraint without being prompted, and finds that prospective recall degrades severely as context grows \citep{zhang2026triggerbench}, while proactive-agent work studies when an agent should volunteer help at all \citep{lu2024proactiveagent}. The failure modes are symmetric and both costly: intrusion when the agent acts on a trigger it should have suppressed, and missed-help when it fails to act on one it should have surfaced.

Trigger state stresses actionability and authority: a trigger is an executable commitment to act under a condition, and the agent needs current authority to fire it, since a scheduled action authorized last week may no longer be permitted today. The gap is that trigger calibration is almost entirely unmeasured as a governance problem; the rare benchmarks score whether the agent can recall a latent constraint, not whether it correctly distinguishes a trigger it is still authorized and scoped to fire from one whose authority has lapsed, leaving the when-to-act decision among the thinnest cells in the corpus.

\subsubsection{Tool and credential state}

Tool and credential state is the agent's authority over external resources: API tokens, OAuth scopes, device-access grants, and the live handles through which the agent reaches the world. It is durable, control-relevant, and emphatically not memory, yet its corruption converts a benign agent into a dangerous one. The threat is best documented from the attack side. Cross-session injection shows that the prompt-injection threat model expands from a single session to the whole system once persistent state, including tool registries and credentials, carries the payload forward \citep{xie2026crosssessioninjection}, and tool drift shows that stored biases accumulate to skew tool selection over time \citep{dabas2026tooldrift}. The recovery side surfaces a subtle hazard named by ACR-Fence: because an agent regenerates slightly different requests after a checkpoint restore, retried tool calls become new external requests that can resurrect already-consumed authority and duplicate irreversible side effects unless those effects are recorded and replayed rather than re-issued \citep{zheng2026acrfence}.

The failure modes are permission drift and stale authorization: a token outlives the consent that justified it, or a scope granted for one operation is reused for another. These are authority and recoverability failures, and they intersect with the most under-served lifecycle stage, rollback. Tool and credential state stresses authority and recoverability above all, and the gap is that the corpus usually treats credentials as configuration rather than governed state with a lifecycle: there is almost no coded work on aging out a token when its authority lapses, on auditing which credential licensed which side effect, or on rolling back tool-originated state once the authority that licensed the original action is gone. ACR-Fence frames the failure, but the corpus has not yet generalized it.

\subsubsection{External commitments and side effects}

External commitments and side effects are not internal state; they are the consequences an always-on agent has written into the world and must keep linked to internal state: a payment issued, a message sent, a file deleted, a record created in an external system. They often have no literal undo, because the agent does not own them, and they are the reason recoverability is more than an internal bookkeeping property. AppWorld is the clearest probe, executing real operations such as payment requests, library deletions, and file creation whose effects persist outside the agent and must be tracked across sessions \citep{trivedi-etal-2024-appworld}. SagaLLM's compensable-transaction model is the most principled response, providing compensation handlers that semantically undo a step when the literal state cannot be restored \citep{chang2025sagallm}, while ACR-Fence shows the failure that arises when side effects are replayed rather than compensated after a restart: duplicated irreversible actions \citep{zheng2026acrfence}.

The failure modes are duplicate side effects and the resurrection of already-consumed authority, both recoverability failures of the highest-consequence kind, because the affected object lives outside the agent's control. External commitments stress recoverability and provenance: every side effect should carry a handle back to the records that justified it, so that a later-discovered bad record yields a bounded set of external effects to compensate, the rollback-traceability invariant made operational. In our coded corpus, only $27$ of $435$ works expose any rollback mechanism at all, and we found no coded work that scores rollback of tool-originated side effects under lapsed authority; this is the cell where thin return-arc coverage has the most direct real-world cost.

\subsection{A comparison across types}

Table~\ref{tab:state-tax} consolidates the taxonomy against characteristic failure modes and the state axes each item stresses. Read down the table, a pattern emerges. Classical memory objects are richly represented, with mature benchmarks and dedicated mechanisms, and their failures, while real, are largely recoverable because the affected state is internal. Object states with external consequences and control-envelope fields are less consistently represented, often by recent work or by attack papers rather than mature mechanisms. The axes column makes the skew precise: authority and recoverability, the two rarest axes in our coded corpus, concentrate in the envelope and in externally consequential object states, while mutability and provenance dominate the well-studied memory core.

\begin{table}[t]
\centering
\small
\setlength{\tabcolsep}{4pt}\renewcommand{\arraystretch}{1.2}
\begin{tabularx}{\linewidth}{p{0.15\linewidth}p{0.23\linewidth}YY}
\toprule
Layer & State item & Typical failure mode & Axes most stressed \\
\midrule
Object state & Episodic traces & Noise, privacy leakage, loss of evidence handles & Recoverability, provenance \\
Object state & Semantic facts & Staleness, contradiction, authority confusion & Mutability, authority \\
Object state & Preferences / user model & Over-personalization, scope creep, drift & Scope, actionability \\
Object state & Procedural skills & Procedure poisoning, unsafe reuse, overfitting & Actionability, authority, recoverability \\
Object state & Workflow / task commitments & False continuity, stale commitments & Mutability, recoverability \\
Object state & Shared / social state & Speaker confusion, propagation, authority conflict & Scope, provenance \\
Object state & Trigger conditions & Intrusion, missed-help & Actionability, authority \\
\midrule
Governance envelope & Authority / permission & Privilege drift, cross-user leakage & Authority, scope \\
Governance envelope & Scope / ownership & Cross-principal use, consent mismatch & Scope, authority \\
Governance envelope & Tool / credential capability & Permission drift, stale authorization & Authority, recoverability \\
Governance envelope & Provenance / audit & Corroboration laundering, rollback gaps & Provenance, recoverability \\
Governance envelope & Retention / deletion / rollback handles & Residual derived state, unbounded repair & Recoverability, provenance \\
Governance envelope & Trust / confidence / temporal validity & Stale trust, unverified promotion, expiry failure & Mutability, authority \\
\midrule
External commitments & Tool/world side effects & Duplicate effects, consumed-authority resurrection & Recoverability, provenance \\
\bottomrule
\end{tabularx}
\caption{The persistent-state taxonomy. Object state names what the agent carries forward. The governance envelope names the fields and capabilities that determine whether an object may influence action and how it can be repaired. External commitments are consequences linked to internal state rather than internal state themselves. The rare axes, authority and recoverability, concentrate in the envelope and in externally consequential objects.}
\label{tab:state-tax}
\end{table}

\subsection{How the types interact, and where governance is thinnest}

The types are not independent, and treating them as separable categories, which the inherited memory taxonomy encourages, is itself a source of failure. They form a derivation order. Episodic traces are consolidated into semantic facts and preferences; successful episodes are distilled into procedural skills; open episodes accrete into the task ledger; every derived record should carry provenance back to the episodes that produced it; and any of these may be shared, may arm a trigger, may exercise a credential, and may ultimately produce an external side effect. A failure at any upstream type propagates downstream along this derivation, which is why the per-type failure modes are not isolable. Provenance-role collapse in the audit type \citep{jin2026provenancerolecollapse} corrupts the authority confusion of semantic facts; memory contagion in shared memory \citep{liu2026memorycontagion} poisons the procedural skills agents distill from a polluted commons; a poisoned preference under expanded scope \citep{xu2026toxicchats} licenses a credential to fire a trigger that produces an irreversible side effect. The thesis that an always-on agent is a single persistent-state system rather than a bundle of memory stores is precisely the claim that these types must be governed jointly, because their failures compose.

The interaction also explains why deletion is the hardest operation and rollback the rarest mechanism. A forget request against a semantic fact must propagate along the entire derivation: to the episodes it was extracted from, the summaries and embeddings that absorbed it, the skills conditioned on it, the shared copies other principals hold, and any external side effect it justified. This is the deletion-propagation invariant, and it is difficult because the types interact; an immutable provenance ledger that makes the derivation auditable \citep{wright2025immutablemem,ouyang2026memlineage} simultaneously makes deletion harder to honor, so the two goals the audit type is supposed to serve, accountability and erasure, are in tension and the corpus has not reconciled them. The transactional and compensation models \citep{chang2025sagallm,zheng2026acrfence} are the closest the literature comes to honoring the recoverability of the downstream types, and they remain confined to a handful of works.

Governance is thinnest, then, where the authority and recoverability axes dominate, where the affected consequence is external or irreversible, and where the lifecycle return arc, forget, audit, and rollback, must run. The classical memory types are better served because their dominant axes, mutability and provenance, are common in the corpus and their failures are often internally recoverable. Governance-envelope fields and external commitments are served by a scattering of mostly recent preprints and by attack papers that document the failure without governing it. The bibliographic skew we measure, only $27$ of $435$ coded works exposing any rollback mechanism and authority being the rarest axis at $72$ of $435$, is the quantitative shadow of this taxonomic skew. The works we coded contain a rich science of the recoverable memory core and a thinner account of the governance envelope around consequential state. That asymmetry, more than any single missing mechanism, is the open problem this taxonomy is built to make visible, and it sets the agenda for the substrate, lifecycle, and governance parts that follow.
